\chapter{Arhitectura sistemului}
\label{ch:arhitectura}

\section{Vedere de ansamblu}

\Cref{fig:arch} prezintă arhitectura de implementare. Toate componentele rulează ca containere
Docker orchestrate cu Docker Compose, într-o rețea virtuală comună.

\begin{figure}[H]
  \centering
  \resizebox{\linewidth}{!}{% Arhitectura distribuită (docker compose).
\begin{tikzpicture}
  \node[comp, minimum width=32mm] (cli) at (0,0) {Clienți / \texttt{loadgen}\\[-1pt]\scriptsize Zipf, closed/open-loop};
  \node[comp, minimum width=32mm, fill=c1!12] (lb) at (0,-1.6) {nginx\\[-1pt]\scriptsize least\_conn, keep-alive};
  \node[comp, minimum width=34mm, minimum height=14mm, fill=c3!10] (a1) at (-2.4,-3.5) {\textbf{app1} (axum/tokio)\\[-1pt]\scriptsize L1 \texttt{memtable}\\[-2pt]\scriptsize single-flight · write-behind};
  \node[comp, minimum width=34mm, minimum height=14mm, fill=c3!10] (a2) at (2.4,-3.5) {\textbf{app2} (axum/tokio)\\[-1pt]\scriptsize L1 \texttt{memtable}\\[-2pt]\scriptsize single-flight · write-behind};
  \begin{scope}[on background layer]
    \node[draw=c3, dashed, rounded corners, fit=(a1)(a2), inner sep=2.5mm] (apps) {};
  \end{scope}
  \node[note, anchor=south west] at (apps.north west) {instanțe web fără stare};
  \draw[arr] (cli) -- (lb);
  \draw[arr] (lb) -- (a1); \draw[arr] (lb) -- (a2);

  % Backend-uri.
  \node[comp, minimum width=15mm, fill=c2!12] (n1) at (-5.6,-6.4) {cache1};
  \node[comp, minimum width=15mm, fill=c2!12] (n2) at (-3.9,-6.4) {cache2};
  \node[comp, minimum width=15mm, fill=c2!12] (n3) at (-2.2,-6.4) {cache3};
  \begin{scope}[on background layer]
    \node[draw=c2, dashed, rounded corners, fit=(n1)(n3), inner sep=2mm] (ring) {};
  \end{scope}
  \node[note, below=1mm of ring] {cluster \texttt{cache-node}\\consistent hashing, TCP};
  \node[comp, minimum width=15mm, fill=c5!10] (rd) at (0,-6.4) {Redis};
  \node[store, fill=c4!10] (pm) at (2.1,-6.4) {PG in-mem\\[-2pt]\scriptsize tmpfs};
  \node[store, minimum width=26mm, minimum height=13mm, fill=c1!10] (db) at (5.3,-6.4) {\textbf{PostgreSQL}\\[-2pt]\scriptsize sursa de adevăr};

  \foreach \t/\c in {ring/c2, rd/c5, pm/c4, db/c1} {
    \draw[arr, \c] (apps.south -| \t.north) -- (\t.north);
  }
  \draw[arr, dashed, c3] (db.east) to[out=60,in=0] node[note, right, pos=0.5] {LISTEN/NOTIFY\\invalidare L1} (apps.east);

  \node[comp, minimum width=26mm, fill=c6!10] (obs) at (5.3,-1.6) {Prometheus\\[-1pt]+ Grafana};
  \draw[arr, dotted] (obs.south) -- node[note, right] {/metrics} (obs.south |- apps.north);
\end{tikzpicture}
}
  \caption{Arhitectura distribuită a sistemului.}
  \label{fig:arch}
\end{figure}

\begin{table}[H]
\centering\small
\begin{tabularx}{\linewidth}{lcX}
\toprule
Componentă & Instanțe & Rol \\
\midrule
\texttt{nginx} & 1 & load balancer HTTP (\texttt{least\_conn}, conexiuni keep-alive către backend) \\
\texttt{webapp} & 2 & API REST (axum); cache L1 \texttt{memtable}; single-flight; write-behind; metrici \\
\texttt{cache-node} & 3 & server TCP peste \texttt{memtable}; partiționare prin consistent hashing \\
PostgreSQL & 1 & baza relațională — sursa de adevăr; trigger \texttt{NOTIFY} \\
PostgreSQL in-memory & 1 & date pe \texttt{tmpfs}, \texttt{fsync=off}; tabela \texttt{UNLOGGED cache\_kv} \\
Redis & 1 & cache de referință (\emph{baseline} industrial), fără persistență \\
Prometheus, Grafana & 1+1 & colectare metrici (2 s) și dashboard provizionat automat \\
\texttt{loadgen} & la cerere & generator de încărcare și benchmark-uri offline \\
\bottomrule
\end{tabularx}
\caption{Componentele sistemului.}
\end{table}

\section{Organizarea codului}

Proiectul este un \emph{workspace} Cargo cu patru crate-uri, cu dependențe aciclice:

\begin{center}
\begin{tikzpicture}[node distance=6mm and 12mm]
  \node[comp, fill=c4!10] (m) {\texttt{memtable}\\[-1pt]\scriptsize bibliotecă};
  \node[comp, fill=c2!12, above right=2mm and 14mm of m] (c) {\texttt{cache-node}\\[-1pt]\scriptsize lib + server};
  \node[comp, fill=c3!12, right=14mm of c] (w) {\texttt{webapp}\\[-1pt]\scriptsize server HTTP};
  \node[comp, fill=c6!10, below right=2mm and 14mm of m] (l) {\texttt{loadgen}\\[-1pt]\scriptsize CLI benchmark};
  \draw[arr] (c) -- (m); \draw[arr] (w) -- (c); \draw[arr] (w.south) to[out=-150,in=10] (m.east); \draw[arr] (l) -- (m);
\end{tikzpicture}
\end{center}

\begin{itemize}
  \item \texttt{memtable} — motorul de tabele in-memory, fără dependențe de rețea; reutilizat de nodurile
        de cache, de aplicație (L1) și de benchmark-uri.
  \item \texttt{cache-node} — protocolul, serverul TCP, clientul cu pool de conexiuni, inelul de
        consistent hashing și clientul de cluster.
  \item \texttt{webapp} — aplicația web: rute, OpenAPI, backend-uri de cache, strategii, metrici, UI.
  \item \texttt{loadgen} — încărcare HTTP, simulare de politici, scalabilitatea motorului.
\end{itemize}

\section{Motorul de tabele in-memory}

\begin{figure}[H]
  \centering
  \resizebox{0.92\linewidth}{!}{% Structura internă a unei tabele memtable.
\begin{tikzpicture}[
  slot/.style={draw, minimum width=13mm, minimum height=7mm, font=\scriptsize, fill=white, align=center},
  sh/.style={draw, rounded corners=2pt, minimum width=9mm, minimum height=6mm, font=\scriptsize, fill=c1!12}]
  \node[comp, minimum width=26mm] (k) {\texttt{get("product:42")}};
  \node[comp, right=8mm of k, minimum width=24mm, fill=c4!10] (h) {hash(key) \& (S$-$1)};
  \draw[arr] (k) -- (h);
  \foreach \i in {0,...,5} { \node[sh] (s\i) at ($(h.south)+(-3.3+1.32*\i,-1.2)$) {\ifnum\i=5 S$-$1\else\ifnum\i=4 …\else\i\fi\fi}; }
  \draw[arr] (h.south) -- (s2.north);
  \node[note, left=1mm of s0] {shard-uri\\(Mutex)};

  % Detaliul unui shard.
  \node[draw, rounded corners, fill=boxbg, minimum width=118mm, minimum height=38mm, below=7mm of s2.south, anchor=north] (box) {};
  \node[font=\scriptsize\bfseries, anchor=north west] at (box.north west) {Shard 2};
  \node[slot, fill=c3!20] (e1) at ($(box.west)+(20mm,2mm)$) {k$_7$\\v, ttl};
  \node[slot, right=6mm of e1] (e2) {k$_{42}$\\v, ttl};
  \node[slot, right=6mm of e2] (e3) {k$_3$\\v, ttl};
  \node[slot, right=6mm of e3, fill=c5!15] (e4) {k$_{19}$\\v, ttl};
  \foreach \a/\b in {e1/e2,e2/e3,e3/e4} {
    \draw[arr, shorten >=1pt] ([yshift=1.2mm]\a.east) -- ([yshift=1.2mm]\b.west);
    \draw[arr, shorten >=1pt] ([yshift=-1.2mm]\b.west) -- ([yshift=-1.2mm]\a.east);
  }
  \node[note, above=1mm of e1] {head (MRU)};
  \node[note, above=1mm of e4] {tail → evacuare};
  \node[draw, font=\scriptsize, align=left, fill=white, anchor=north west] (map) at ($(box.south west)+(3mm,13mm)$) {\texttt{HashMap<key, slot>}\\ căutare O(1)};
  \node[draw, font=\scriptsize, align=left, fill=white, right=5mm of map] (tags) {\texttt{tags: HashMap<tag, \{slot\}>}\\ invalidare pe grup};
  \node[draw, font=\scriptsize, align=left, fill=white, right=5mm of tags] (free) {\texttt{free: Vec<slot>}\\ reutilizare sloturi};
  \draw[arr, densely dotted] (map.north) to[out=90,in=-90] (e2.south);
  \node[note, below=1mm of box] {LRU: la acces, nodul se mută în cap · FIFO: nu se mută · CLOCK: bit de referință, „a doua șansă" la coadă};
\end{tikzpicture}
}
  \caption{Structura unei tabele \texttt{memtable}: shard-uri independente, fiecare cu un hash map,
           o listă dublu înlănțuită intrusivă peste un slab și un index de tag-uri.}
  \label{fig:memtable}
\end{figure}

O tabelă (\cref{fig:memtable}) este împărțită în $S$ shard-uri ($S$ putere a lui 2). Cheia este
direcționată către shard prin $\text{hash}(k) \mathbin{\&} (S-1)$, iar fiecare shard este protejat de
propriul mutex. În interiorul shard-ului:

\begin{itemize}
  \item intrările sunt stocate într-un \emph{slab} (\texttt{Vec<Slot>}); legăturile listei sunt indici
        \texttt{u32}, nu pointeri — fără \texttt{unsafe}, cu localitate bună în cache-ul procesorului;
  \item un \texttt{HashMap<cheie, index>} oferă căutare O(1);
  \item lista dublu înlănțuită ordonează intrările pentru evacuare: capul este cea mai recentă, coada este
        candidatul; toate operațiile (acces, inserare, ștergere, evacuare) sunt O(1);
  \item indexul de tag-uri \texttt{HashMap<tag, \{index\}>} permite invalidarea unui grup de chei (de
        exemplu toate clasamentele unei categorii) fără a parcurge tabela;
  \item sloturile eliberate sunt reținute într-o listă liberă și refolosite (fără realocări).
\end{itemize}

Valorile sunt \texttt{bytes::Bytes} — buffere imutabile cu numărare de referințe, deci un hit nu copiază
datele: același buffer ajunge direct în corpul răspunsului HTTP.

\section{Cache-ul distribuit}

\begin{figure}[H]
  \centering
  \begin{subfigure}{0.42\linewidth}\centering
    \resizebox{\linewidth}{!}{% Inelul de consistent hashing cu noduri virtuale.
\begin{tikzpicture}
  \def\R{2.2}
  \draw[thick, gray!60] (0,0) circle (\R);
  \foreach \ang/\col/\lab in {10/c1/A,40/c2/B,75/c3/C,110/c1/A,150/c3/C,185/c2/B,220/c1/A,255/c2/B,290/c3/C,330/c2/B} {
    \fill[\col] (\ang:\R) circle (3.2pt);
    \node[font=\scriptsize, text=\col] at (\ang:\R+0.35) {\lab};
  }
  \foreach \ang/\k in {25/k_1,95/k_2,200/k_3,305/k_4} {
    \draw[fill=white, draw=black] (\ang:\R) circle (2.2pt);
    \node[font=\scriptsize] at (\ang:\R-0.42) {$\k$};
  }
  \draw[-{Stealth}, thick, black!60] (27:\R-0.15) arc (27:38:\R-0.15);
  \draw[-{Stealth}, thick, black!60] (97:\R-0.15) arc (97:108:\R-0.15);
  \node[note, align=center] at (0,0) {$h(\text{cheie}) \mapsto$ primul\\nod virtual în sens orar\\[2pt] 160 noduri virtuale / nod};
\end{tikzpicture}
}
    \caption{Inelul de consistent hashing.}
    \label{fig:ring}
  \end{subfigure}\hfill
  \begin{subfigure}{0.54\linewidth}\centering\small
    \begin{tabular}{@{}ll@{}}
      \toprule
      Comandă & Răspuns \\
      \midrule
      \texttt{GET k} & \texttt{\$len\textbackslash n bytes} / \texttt{\$-1} \\
      \texttt{SET k ttl len [tags]} & \texttt{+OK} \\
      \texttt{DEL k} & \texttt{:1} / \texttt{:0} \\
      \texttt{TAG t} & \texttt{:n} \\
      \texttt{STATS} & \texttt{\$len\textbackslash n json} \\
      \texttt{CONFIG cap pol shards} & \texttt{+OK} \\
      \bottomrule
    \end{tabular}
    \caption{Protocolul nodului (extras).}
  \end{subfigure}
  \caption{Partiționarea și protocolul cache-ului distribuit.}
\end{figure}

Nodurile nu comunică între ele: partiționarea este realizată în client. Fiecare nod fizic este plasat
de 160 de ori pe un inel de $2^{64}$ poziții, folosind o funcție de dispersie stabilă (FNV-1a urmată de
un finalizator \emph{splitmix64}). Căutarea nodului responsabil se face prin căutare binară în vectorul
sortat al punctelor, în $O(\log(160n))$. Testele unitare verifică echilibrul (fiecare din 3 noduri
primește între 27\% și 40\% din chei) și faptul că adăugarea unui al patrulea nod mută
$\approx 25\%$ din chei.

Protocolul este textual pentru antete și binar pentru valori (prefixate cu lungimea), deci poate
transporta orice conținut. Serverul trimite răspunsurile (\emph{flush}) doar când bufferul de intrare
s-a golit, ceea ce permite \emph{pipelining}: un client poate trimite mai multe comenzi fără să aștepte
fiecare răspuns.

\section{Fluxul unei cereri de citire}

\begin{figure}[H]
  \centering
  \resizebox{0.95\linewidth}{!}{% Diagrama de secvență: citire cache-aside cu single-flight.
\begin{tikzpicture}[font=\small]
  \foreach \x/\n/\t in {0/c1/Client 1,2.6/c2/Client 2,5.8/app/Aplicație,8.8/cache/Cache,11.6/db/PostgreSQL} {
    \node[comp, minimum width=20mm] (\n) at (\x,0) {\t};
    \draw[gray!60, dashed] (\n.south) -- ++(0,-7.4);
  }
  \newcommand{\msg}[5][]{\draw[arr,#1] (#2,#4) -- node[above, font=\scriptsize]{#5} (#3,#4);}
  \msg{0}{5.8}{-1.0}{GET /products/42}
  \msg{2.6}{5.8}{-1.5}{GET /products/42}
  \msg{5.8}{8.8}{-2.1}{get(k)}
  \msg[dashed]{8.8}{5.8}{-2.6}{miss}
  \node[draw, fill=c4!10, font=\scriptsize, align=center] at (5.8,-3.35) {single-flight: o singură\\interogare per cheie};
  \msg{5.8}{11.6}{-4.2}{SELECT … WHERE id = \$1}
  \msg[dashed]{11.6}{5.8}{-4.8}{rând}
  \msg{5.8}{8.8}{-5.5}{set(k, v, ttl)}
  \msg[dashed]{5.8}{0}{-6.3}{200, X-Cache: MISS}
  \msg[dashed]{5.8}{2.6}{-6.9}{200 (coalescat)}
\end{tikzpicture}
}
  \caption{Citire cache-aside cu single-flight: două cereri concurente pentru aceeași cheie produc o
           singură interogare SQL.}
  \label{fig:readpath}
\end{figure}

La un miss, aplicația nu interoghează imediat baza de date: verifică într-o tabelă de „zboruri în
curs" (\texttt{HashMap<cheie, Arc<OnceCell>>}) dacă o altă cerere încarcă deja aceeași cheie. Dacă da,
așteaptă rezultatul acesteia (\cref{fig:readpath}). Mecanismul elimină \emph{cache stampede} la pornire
sau după invalidarea unei chei fierbinți (experimentul E7).

\section{Strategiile de scriere}

\begin{figure}[H]
  \centering
  \resizebox{\linewidth}{!}{% Cele trei strategii de scriere.
\begin{tikzpicture}[font=\small, node distance=4mm]
  \foreach \col/\x/\title/\a/\b/\c in {
      c1/0/Invalidate (cache-aside)/1. UPDATE în DB/2. DEL cheie din cache/3. următoarea citire: miss → DB,
      c3/5.2/Write-through/1. UPDATE în DB/2. SET valoarea nouă în cache/3. citirile rămân hit,
      c2/10.4/Write-behind/1. SET în cache; răspuns 202/2. coadă → lot (≤1024)/3. un UPDATE … unnest()} {
    \node[draw=\col, very thick, rounded corners, minimum width=48mm, minimum height=27mm, fill=\col!6] (b) at (\x,0) {};
    \node[font=\small\bfseries, text=\col, anchor=north] at (b.north) {\title};
    \node[font=\scriptsize, align=left, anchor=north west] at ($(b.north west)+(2mm,-7mm)$) {\a\\[2pt]\b\\[2pt]\c};
  }
  \node[note, text width=46mm] at (0,-2.0) {+ simplu, corect\\− miss după fiecare scriere};
  \node[note, text width=46mm] at (5.2,-2.0) {+ hit ratio păstrat\\− latență de scriere mai mare};
  \node[note, text width=46mm] at (10.4,-2.0) {+ scrieri rapide, grupate\\− consistență eventuală, risc la crash};
\end{tikzpicture}
}
  \caption{Strategiile de scriere implementate.}
  \label{fig:writes}
\end{figure}

Pentru write-behind (\cref{fig:writes}, dreapta), modificările sunt distribuite pe 4 benzi paralele după
\texttt{id mod 4}. Toate modificările unui produs trec prin aceeași bandă, deci ordinea lor se păstrează.
Fiecare bandă este un canal mărginit (25.000 de elemente), care oferă \emph{backpressure}: dacă baza de
date nu ține pasul, cererile de scriere așteaptă în loc să acumuleze o coadă nelimitată. Workerul fiecărei
benzi preia până la 1.024 de modificări deodată, le combină per produs (ultima scriere câștigă, pe câmp) și
le aplică printr-un singur \texttt{UPDATE \ldots{} FROM unnest(\ldots)}. După flush, invalidarea
clasamentelor categoriilor afectate rulează concurent, într-un task separat.
La oprirea aplicației (SIGTERM), canalul este închis și golit complet înainte de terminarea procesului.

\section{Coerența cache-urilor locale}

Fiecare instanță are propriul L1. Un trigger \texttt{AFTER UPDATE} pe tabela \texttt{products} execută
\texttt{pg\_notify} pe canalul \texttt{cache\_invalidation}, cu mesajul \texttt{id:categorie}; fiecare instanță ține o conexiune dedicată
cu \texttt{LISTEN} și, la primirea notificării, șterge cheia produsului și tag-ul categoriei din L1.
Avantajul major: și modificările făcute direct în baza de date (scripturi, alte aplicații) invalidează
cache-ul. Pentru că \texttt{NOTIFY} serializează commit-urile (\cref{sec:optimizari}), aplicația îl
suprimă (\texttt{app.notify = off}) când backend-ul activ nu are nivel local.

\section{Reziliență}

\begin{itemize}
  \item \textbf{Degradare grațioasă}: orice eroare a backend-ului de cache este contorizată și cererea
        continuă direct pe baza de date.
  \item \textbf{Conexiuni}: clientul de nod reîncearcă o dată cu o conexiune nouă dacă o conexiune din pool
        s-a închis (nod repornit); conectarea la baze de date se reîncearcă la pornire.
  \item \textbf{Ascultătorul NOTIFY} se reconectează automat.
  \item \textbf{Write-behind}: reîncercări cu \emph{backoff} exponențial (8 încercări) și golirea cozii la oprire.
\end{itemize}
